% ── Περίληψη (Ελληνική) ──────────────────────────────────────

\noindent
Η αυξανόμενη ζήτηση για έξυπνη παρακολούθηση υγείας σε περιβάλλοντα
Υποβοηθούμενης Διαβίωσης (Ambient Assisted Living, AAL) καθιστά την
Αναγνώριση Δραστηριότητας Ανθρώπου (Human Activity Recognition, HAR)
κρίσιμο ερευνητικό πεδίο. Οι κλασικές κεντρικές προσεγγίσεις μηχανικής
μάθησης απαιτούν τη μεταφορά ευαίσθητων δεδομένων αισθητήρων κινητών
συσκευών σε κεντρικούς εξυπηρετητές, εγείροντας σοβαρές ανησυχίες για
την ιδιωτικότητα των χρηστών.

Η παρούσα διπλωματική εργασία προτείνει ένα σύστημα Ομοσπονδιακής
Μάθησης (Federated Learning, FL) για HAR σε περιβάλλοντα AAL, το οποίο
διατηρεί τα δεδομένα τοπικά σε κάθε συσκευή και μεταδίδει μόνο
ενημερώσεις παραμέτρων μοντέλου στον κεντρικό εξυπηρετητή. Θεμελιώδης
αδυναμία των υπαρχόντων συστημάτων FL είναι η έλλειψη μηχανισμών που
να παρακινούν ορθολογικούς χρήστες να συμμετέχουν ενεργά και να
συνεισφέρουν δεδομένα υψηλής ποιότητας. Για την αντιμετώπιση αυτής
της αδυναμίας, σχεδιάζεται και υλοποιείται μηχανισμός κινήτρων βασισμένος
στο Παίγνιο Stackelberg της Θεωρίας Παιγνίων, στον οποίο ο κεντρικός
εξυπηρετητής ενεργεί ως ηγέτης (leader) καθορίζοντας τη δομή ανταμοιβών,
ενώ οι πελάτες (clients) ενεργούν ως ακόλουθοι (followers) βελτιστοποιώντας
την προσπάθειά τους.

Επιπλέον, για την αξιολόγηση της συνεισφοράς κάθε συμμετέχοντος πελάτη,
αναπτύσσεται σύστημα βασισμένο στα Shapley Values της Συνεργατικής Θεωρίας
Παιγνίων, εξασφαλίζοντας δίκαιη κατανομή ανταμοιβών ανάλογη με την οριακή
συνεισφορά κάθε πελάτη στην αναγνωριστική ικανότητα του καθολικού μοντέλου.
Η προτεινόμενη αρχιτεκτονική υλοποιείται με μονοδιάστατο Συνελικτικό
Νευρωνικό Δίκτυο (1D-CNN) ειδικά σχεδιασμένο για ανάλυση χρονοσειρών
αισθητήρων, με υλοποίηση του πρωτοκόλλου ομοσπονδιακής εκπαίδευσης
(FedAvg) εξ ολοκλήρου σε PyTorch.

Η αξιολόγηση διεξάγεται σε δύο σύνολα δεδομένων: το καθιερωμένο UCI
Human Activity Recognition dataset υπό ελεγχόμενες συνθήκες, και το
AAL Smartphone Dataset για HAR σε AAL υπό
πραγματικές συνθήκες με ετερογενείς, Non-IID κατανομές δεδομένων μεταξύ
22 χρηστών-clients.
Η αρχική αξιολόγηση επιβεβαιώνει τη λειτουργικότητα του συστήματος:
ομοσπονδιακή ακρίβεια $76{,}8\%$ έναντι $84{,}8\%$ κεντρικοποιημένης,
ποσοτικοποιώντας το κόστος της Non-IID ετερογένειας. Η ποσοτική σύγκριση
των μετρικών συνεισφοράς (Leave-One-Out έναντι Shapley) για την
οριστικοποίηση του μηχανισμού κινήτρων αποτελεί το επόμενο στάδιο.

\vspace{1.5em}
\noindent
\textbf{Λέξεις-κλειδιά:}
Ομοσπονδιακή Μάθηση,
Αναγνώριση Δραστηριότητας Ανθρώπου (HAR),
Υποβοηθούμενη Διαβίωση (AAL),
Ψηφιακό Δίδυμο,
Θεωρία Παιγνίων,
Παίγνιο Stackelberg,
Shapley Values,
Μηχανισμοί Κινήτρων.
